\documentclass[10pt,a4paper]{article}
\usepackage[utf8]{inputenc}
\usepackage[T1]{fontenc}
\usepackage[italian]{babel}
\usepackage{amsmath}
\usepackage{amsthm}
\usepackage{amsfonts}
\usepackage{amssymb}
\usepackage{graphicx}
\usepackage{listings}
\usepackage{color}
\usepackage{hyperref}
\usepackage{a4wide}
\definecolor{mygreen}{RGB}{28,172,0} % color values Red, Green, Blue
\definecolor{mylilas}{RGB}{170,55,241}
\setlength{\parindent}{0pt}
\theoremstyle{definition}% default
\newtheorem{algo}{Algorithm}
\author{Luca Formaggia}
\date{AY 2021-22}
\title{A Note on Synthetic Division and Polynomial Zeros}
\begin{document}
  \lstset{language=c++,%
     %basicstyle=\color{red},
     breaklines=true,%
     morekeywords={matlab2tikz},
     keywordstyle=\color{blue},%
     morekeywords=[2]{1}, keywordstyle=[2]{\color{black}},
     identifierstyle=\color{black},%
     stringstyle=\color{mylilas},
     commentstyle=\color{mygreen},%
     showstringspaces=false,%without this there will be a symbol in the places where there is a space
     numbers=left,%
     numberstyle={\tiny \color{black}},% size of the numbers
     numbersep=9pt, % this defines how far the numbers are from the text
     emph=[1]{for,end,break},emphstyle=[1]\color{red}, %some words to emphasise
     %emph=[2]{word1,word2}, emphstyle=[2]{style},    
 }
 
 \maketitle
 \section{Synthetic Division and Horner's Algorithm}
 Let
 \[
 p_m(x)=\sum_{i=0}^{m} a_i x^i,
 \]
be a (real or complex) polynomial of degree at most $m$, expressed in terms of the monomial basis.
  The \emph{synthetic division} algorithm, or \emph{Horner's algorithm}, reads

\begin{algo}\label{alg:horner}
\emph{Computes $p(z)$}
\begin{lstlisting}
auto horner(a,z)
    b[m-1]=a[m];
    for k=m-2,...,0 
	   b[k]+=b[k+1]*z + a[k+1]; 
    return {b[0]*z + a[0],b}; 
\end{lstlisting}
\end{algo}
This version of the algorithm, besides computing $p_m(z)$,
  returns the coefficients $b_j$, $j=0,\ldots m-1$, of the so-called \emph{associated polynomial}
\begin{equation}
\label{eq:associated}
q_{m-1}(x;z)=\sum_{j=0}^{m-1}b_j x^j.
\end{equation}
It is a polynomial of degree at most $m$, and the $z$ in $q_{m-1}(x;z)$ indicates
  that the coefficients $b_j$ are obtained by applying synthetic division at $x=z$ (see Algorithm~\ref{alg:horner}). Indeed, the values of the coefficients $b_j$ depend on the chosen point $z$.

We have the following \emph{important result}: the polynomial $q_{m-1}(x;z)$ is the quotient
  obtained by dividing $p_m$ by $x-z$. This means that
\begin{equation}
\label{eq:quotient}
\boxed{p_m(x)=q_{m-1}(x;z)(x-z) + p_m(z),}
\end{equation}
and, consequently,
\begin{equation}
\label{eq:quotientder}
\boxed{\frac{d}{dx}p_m(x)=\frac{d}{dx}q_{m-1}(x;z)(x-z) + q_{m-1}(x;z),}
\end{equation}
hence,
\begin{equation}
\label{eq:quotientdertwo}
\boxed{\frac{d}{dx}p_m(z)=q_{m-1}(z;z).}
\end{equation}

Therefore, synthetic division can be used to compute not only the value at a point $z$, but also the derivative at that point and the quotient obtained by dividing by the factor $x-z$. It provides the basic ingredients of the \textbf{Newton-Horner} algorithm for computing the zeros of a polynomial. But first we need to discuss deflation.
\subsection{Deflation}
Let $\{\alpha_j,\ j=1,\ldots,n\}$ be the zeros of the polynomial $p_n$. Then, $\{\alpha_j,\ j=1,\ldots,n-1\}$ are the zeros of 
the polynomial
\begin{equation}
p_{n-1}=p_n/(x-\alpha_n).
\end{equation} 
The proof is trivial.
Moreover, $p_{n-1}(x)=q_{m-1}(x;\alpha_n)$, according to~\eqref{eq:quotient}, since $p_n(\alpha_n)=0$. This allows us to define a very efficient algorithm for 
finding all the zeros of a polynomial.
\subsection{Newton-Horner with deflation}
The algorithm implemented here is Newton-Horner with deflation:

\begin{algo}
    Given the polynomial $p_n$, the initial point $x_0$, and the tolerances $tolr$ (tolerance on the residual) and $told$
    (tolerance on the iteration step),
    do the following:
    \begin{description}
        \item[] Set $q=p_n$.
        \item[] For $k=1,n$:
        \begin{enumerate}
            \item Perform the Newton iteration: for $n=0,1,\ldots$
            \begin{enumerate}
                \item Compute $q(x_n)$ and $q^\prime(x_n)$ using synthetic division;
                \item Compute $x_{n+1}=x_{n}-q^\prime(x_{n})^{-1}q(x_n)$;
                \item If $|q(x_{n+1})|<tol$ and $|x_{n+1}-x_{n}|<told$, terminate the iteration and set
            $\alpha_k=x_{n+1}$ (the $k$-th zero);
            \end{enumerate}
            \item Set $q=q/(x-\alpha_k)$ (deflation) using the associated polynomial returned by the last synthetic division;
            \item Set $x_0=\overline{\alpha_k}$ (to facilitate the search for conjugate zeros);
        \end{enumerate}
    \end{description}
\end{algo}

This algorithm exploits the fact that synthetic division (Horner's algorithm) provides the value of both the polynomial and its derivative at a point and, once a zero has been found, also the coefficients of the deflated polynomial. 
\subsection{Some comments}
The Newton-Horner algorithm is very elegant and relatively efficient. However:
\begin{itemize}
    \item We have global convergence (i.e. for all $x_0$) only for polynomials with real coefficients, and only under some conditions on the zeros; see the specialized literature for more details. Therefore, in general, convergence of the basic algorithm depends on the choice of $x_0$. You may implement techniques such as backtracking or randomization to improve convergence properties. In all the tests I carried out, however, the basic method worked well.
    \item Deflation is only approximate. Indeed, we divide not by the exact zero, but by an \emph{approximate zero}. This means that errors may accumulate and reduce the accuracy of the last computed zeros: a small value of the \emph{approximate} deflated polynomial may not imply a small value of the original polynomial. Consequently, the algorithm may terminate at a point far from the true zero. 
    However, this is mainly an issue when computing all the zeros of a high-degree polynomial, especially the last ones. Normally, if the tolerances are reasonably tight, there is no need for concern. 
    \item In polynomials with real coefficients, complex roots come in conjugate pairs. Can we deflate directly by $(x-\alpha)(x-\overline{\alpha})$?
    Indeed we can; this is a modification of the basic algorithm that I have not implemented here.
    \item You must avoid division by zero in $q^\prime(x_n)^{-1}p(x_n)$.
    This is a delicate issue. One possible workaround is to use a small number in the denominator. For instance, use $p(x_n)/d$, with $d=\epsilon \ll 1$ if $|q^\prime(x_n)|=0$, and $d=q^\prime(x_n)$ otherwise. 
    \item Working with complex numbers makes the algorithm more general. However, if $x_0$ is purely real (i.e. it has zero imaginary part), then the algorithm will find only real zeros. Since, in general, we do not know how many real zeros there are, the algorithm may fail to find a zero simply because no additional real zeros exist. Therefore, we need to modify the algorithm to detect non-convergence. Here I decided to keep things simple, but this means that you should always start with an 
    $x_0$ with non-zero imaginary part, unless you already know that \emph{all} zeros are real or, at least, that you know how many real zeros there are, since the given algorithm allows one to search only for a prescribed number of zeros.
    \item When is a complex number actually real? The answer seems straightforward: when the imaginary part is zero. This is true in the ideal world of exact complex arithmetic, but not in the real world of \emph{floating-point complex numbers}. Roundoff errors
    often make the imaginary part very small, but not exactly zero. A rule of thumb is to treat the imaginary part of the approximate root found by the algorithm as zero if its absolute value is smaller than the given tolerances. Of course, if you know a priori that the root is real, you will consider only the real part of the computed approximation.
\end{itemize}


\section{Derivatives}
By leveraging Equations~\eqref{eq:quotientder} and~\eqref{eq:quotient}, we can also compute higher derivatives at the point $z$. Let us first look at the second derivative:
\begin{multline}\label{eq:secderexp}
\frac{d^2p_m}{dx^2}(x)=\frac{d^2q_{m-1}}{dx^2}(x;z)(x-z) + \frac{dq_{m-1}}{dx} (x;z) + \frac{dq_{m-1}}{dx} (x;z)=\\ \frac{d^2}{dx^2}q_{m-1}(x;z)(x-z) + 2 \frac{d}{dx} q_{m-1}(x;z),
\end{multline}
Thus,
\begin{equation}\label{eq:secder}
\frac{d^2p_m}{dx^2}(z)=2 \frac{d}{dx} q_{m-1}(z;z),
\end{equation}

Using~\eqref{eq:quotientder},
\begin{equation}\label{eq:firstderass}
\frac{d}{dx}q_{m-1}(x;z)=\frac{d}{dx}q_{m-2}(x;z)(x-z)+ q_{m-2}(x;z),
\end{equation}
which, combined with~\eqref{eq:secder}, gives
\begin{equation}\label{eq:secderivative}
\frac{d^2p_m}{dx^2}(z)=2 q_{m-2}(z;z).
\end{equation}
So the second derivative at $z$ is twice the associated polynomial of $q_{m-1}(x;z)$ evaluated at the same point. Therefore, it can be computed with just two applications of synthetic division.
What about the third derivative? From~\eqref{eq:secderexp} we have
\begin{multline}\label{eq:secderexptwo}
\frac{d^3p_m}{dx^3}(x)=\frac{d^3}{dx^3}q_{m-1}(x;z)(x-z) +\frac{d^2}{dx^2}q_{m-1}(x;z) + 2 \frac{d^2}{dx^2} q_{m-1}(x;z)=\\
\frac{d^3}{dx^3}q_{m-1}(x;z)(x-z) +3\frac{d^2}{dx^2}q_{m-1}(x;z).
\end{multline}
Using~\eqref{eq:firstderass}
\begin{equation}
\frac{d^2}{dx^2}q_{m-1}(x;z)=\frac{d^2}{dx^2}q_{m-2}(x;z)(x-z)+2\frac{d}{dx}q_{m-2}(x;z),
\end{equation}
and thus
\begin{equation}\label{eq:thirdder}
\frac{d^3p_m}{dx^3}(x)=
\frac{d^3}{dx^3}q_{m-1}(x;z)(x-z) +3\frac{d^2}{dx^2}q_{m-2}(x;z)(x-z)^2 + 6\frac{d}{dx}q_{m-2}(x;z),
\end{equation}
from which, exploiting again~\eqref{eq:quotientder}, we have
\begin{equation}\label{eq:thirdderz}
\frac{d^3}{dx^3}p_m(z)=6 q_{m-3}(z;z).
\end{equation}

In general, for any $n\le m$,
\begin{equation}
\boxed{\frac{d^n}{dx^n}p_m(z)=n! q_{m-n}(z;z).}
\end{equation}

Therefore, all derivatives at a given point $z$ may be computed by repeated application of synthetic division.


\section{The code in \texttt{polyHolder.hpp}}
The code in \texttt{polyHolder.hpp} contains:
\begin{itemize}
    \item The function \lstinline|polyEval()|, which uses the basic Horner algorithm to compute the value of a polynomial at a point;
    \item The class \lstinline|PolyHolder|, which is a helper class for the Newton-Horner algorithm and also allows one to compute derivatives of a polynomial, given its coefficients, at a given point;
    \item The function \lstinline|polyRoots()|, which computes the roots of a polynomial using the basic \emph{Newton-Horner} procedure.
\end{itemize}

One may argue that I am not fully respecting the \emph{single-responsibility principle} here: the class \texttt{PolyHolder} should provide only synthetic division. The computation of derivatives should instead be the responsibility of a function, or of another class, that uses \texttt{PolyHolder}. That objection is reasonable, but the world is not perfect, and \emph{principles} should sometimes be applied with a grain of salt.


\end{document}
